\cotefile{156-7}
\scanpage{2}
\mspage{1}
\margin{\tikz[baseline=(h.base)]{\node[inner sep=2pt] (h) {\unsure{GF|CF} VII}; \draw (h.north west)--(h.south west)--(h.south east);}}
\begin{center}Analysis Situs (troisième mouture)\end{center}
% slanted
\margin{Cf surtout p. 23 -- 50-58, 59 \tikz[baseline=(q.base)]{\node[inner sep=2pt] (q) {110-113}; \draw[->] (q.north)..controls (q.north east) and (q.north east)..(q.east)..controls (q.south east) and (q.south east)..(q.south)..controls (q.south west) and (q.south west)..(q.west)..controls (q.north west) and (q.north west)..(q.north);}}

\uline{23 juin}

Je récapitule les données et axiomes d'une ``algèbre des figures'', ou ``\uline{atelier}''. \margin{--}Il s'agit pour l'instant d'axiomes purement ``algébriques'' dominant le formalisme algébrique des figures, à l'exclusion pour le moment des propriétés topologiques (densité, connexité, acyclicité) ou \inserted{\uline{proprement}} géométriques (régularité topologique, boules et sphères topologiques d'une algèbre des figures), voire les ``propriété combinatoires fines'', \unsure{autour} de la notion de simplexe.

Parmi les axiomes, je distingue les axiomes ``sine qua non'' At~1 à At~7, qui me paraissent les plus essentielles pour fonder le sens géométrico-topologique du formalisme envisagé. Les autres axiomes seront considérés comme plus ou moins ``facultatifs'' au niveau du formalisme algébrique. Ils \unsure{pourront|joueront} sans \unsure{doute} un rôle crucial, dès qu'il s'agira d'un travail sur piéces dans le
